% English translation of questions/5782/semA-special/q1.tex (metadata is taken from the Hebrew file).

\begin{question}
For arbitrary parameters $a,b\in\R$ define a function by
\[ f(x)=\begin{cases} e^{\frac1x} & x<0,\\ ax+b-6 & 0\le x\le 9,\\ \frac{x^2-8x-9}{3-\sqrt{x}} & x>9 \end{cases} \]
For which values of the parameters $a,b$ is the function continuous on $\R$? Justify.
\end{question}

\begin{hint}
On each of the open intervals, $f$ is a composition/quotient of continuous functions. It remains to check the junction points $x=0$ and $x=9$.
\end{hint}

\begin{hint}
At $x=9$: $x^2-8x-9=(x-9)(x+1)$ and $x-9=(\sqrt x-3)(\sqrt x+3)$.
\end{hint}

\begin{solution}
\textbf{Continuity inside the intervals.} On $(-\infty,0)$ the function $e^{1/x}$ is continuous (a composition of continuous functions). On $(0,9)$ $f$ is a polynomial. On $(9,\infty)$ $f$ is a quotient of continuous functions whose denominator $3-\sqrt x<0$ does not vanish. Hence $f$ is continuous at every $x\neq0,9$ for all $a,b$, and we only need to check $x=0$ and $x=9$.

\textbf{The point $x=0$.} $f(0)=b-6$, and the limit from the right is $\lim_{x\to0^+}(ax+b-6)=b-6$. From the left: as $x\to0^-$ we have $\frac1x\to-\infty$, and hence
\[ \lim_{x\to0^-}e^{1/x}=0. \]
Continuity at $0$ is equivalent to $b-6=0$, i.e. $b=6$.

\textbf{The point $x=9$.} $f(9)=9a+b-6$, and this is also the limit from the left. From the right, for $x>9$:
\[ \frac{x^2-8x-9}{3-\sqrt x}=\frac{(x-9)(x+1)}{3-\sqrt x}=\frac{(\sqrt x-3)(\sqrt x+3)(x+1)}{-(\sqrt x-3)}=-(\sqrt x+3)(x+1), \]
and hence
\[ \lim_{x\to9^+}f(x)=-(3+3)(9+1)=-60. \]
Continuity at $9$ is equivalent to $9a+b-6=-60$.

\textbf{Solving the system.} From $b=6$ we get $9a=-60$, i.e. $a=-\frac{20}{3}$.

\textbf{Answer:} $f$ is continuous on all of $\R$ if and only if $\boxed{a=-\tfrac{20}{3},\ b=6}$.
\end{solution}
